\section{2015 Algebra \& Number Theory}

\subsection*{Individual}
\begin{exercise}
Let $G$ be a finite abelian group with a bilinear, alternating, and non-degenerate pairing $\ell: G \times G \to \mathbb{Q}/\mathbb{Z}$.
Show that $G \cong H \oplus H$ for some $H$ such that $\ell|_H = 0$.
\end{exercise}

\begin{solution}
\textbf{Motivation:} This is the analogue of the classification of symplectic forms for finite abelian groups. Such pairings are called \textit{symplectic pairings} or \textit{dual pairings}.

(1) \textbf{Image of $\ell_a$:} Let $o(a) = M$. Then $M \ell(a, b) = \ell(Ma, b) = 0 \pmod{\mathbb{Z}}$, so $\text{im}(\ell_a) \subseteq \frac{1}{M}\mathbb{Z}/\mathbb{Z}$.
By non-degeneracy, the map $G \to \widehat{G}$ given by $a \mapsto \ell_a$ is an isomorphism. The order of $\ell_a$ in the dual group is $M$. Since the dual group is isomorphic to $G$, and a cyclic subgroup of order $M$ in $\mathbb{Q}/\mathbb{Z}$ is unique (namely $\frac{1}{M}\mathbb{Z}/\mathbb{Z}$), the image must be exactly that.

(2) \textbf{Hyperbolic Pair:} Pick $a \in G$ of maximal order $M$. Since $\text{im}(\ell_a) = \frac{1}{M}\mathbb{Z}/\mathbb{Z}$, there exists $b$ such that $\ell(a, b) = 1/M$.
Then $o(b) \geq M$ because its image contains an element of order $M$. Since $M$ is maximal, $o(b) = M$.

(3) \textbf{Direct Sum:} Let $S = \langle a, b \rangle$.
If $xa+yb \in S^\perp$, then $\ell(xa+yb, a) = -y/M = 0 \implies M|y$. Similarly $M|x$. Thus $S \cap S^\perp = \{0\}$.
By a counting argument (or explicit projection), $G = S \oplus S^\perp$.
$S \cong \mathbb{Z}/M \oplus \mathbb{Z}/M$ and $\ell$ restricted to $S$ is the standard "symplectic" block.

(4) \textbf{Induction:} Since $S^\perp$ is smaller and $\ell$ restricted to it is still non-degenerate and alternating, we repeat the process.
Eventually $G \cong \bigoplus (\mathbb{Z}/M_i \oplus \mathbb{Z}/M_i)$. Letting $H = \bigoplus \mathbb{Z}/M_i$ gives the result.
\end{solution}

\begin{exercise}
Invariant theory of $O_n(\mathbb{C})$. $F_i = \{ f \in \mathbb{C}(M_{n,k}) : f(gx) = \det(g)^i f(x) \}$.
\end{exercise}

\begin{solution}
(1) $\det Q(x) \neq 0 \iff \dim V_x = k$ and $V_x \cap V_x^\perp = 0$.
$Q(x) = x^t x$. This is the Gram matrix of the columns.
If $\det Q \neq 0$, the columns are linearly independent (so $\dim V_x = k$).
Non-degeneracy follows because the restriction of the Euclidean form to $V_x$ has matrix $Q(x)$.

(2) $F_0$ is the field of $O_n$-invariant rational functions.
\textbf{First Fundamental Theorem for $O_n$:} The ring of invariant polynomials is generated by $Q(x)_{ij} = v_i \cdot v_j$.
Thus $F_0$ is generated by the entries of $Q(x)$.

(3) If $k < n$, $V_x^\perp \neq 0$. Pick $w \in V_x^\perp$ with $w \cdot w \neq 0$.
The reflection $s_w$ fixes $V_x$ and has $\det s_w = -1$.
Then $f(x) = f(s_w x) = \det(s_w) f(x) = -f(x)$, so $f(x) = 0$.
Since this holds for generic $x$, $f \equiv 0$.

(4) If $k \geq n$, let $\delta(x) = \det(v_1, \dots, v_n)$.
Then $\delta(gx) = \det(g) \delta(x)$, so $\delta \in F_1$.
For any $h \in F_1$, $h/\delta$ is invariant ($F_0$), so $F_1 = F_0 \cdot \delta$.
\end{solution}

\begin{exercise}
$x^3+x+1=0$. Show $(31/p)=-1 \implies$ solvable mod $p$.
\end{exercise}

\begin{solution}
(1) $f(x) = x^3+x+1$ is irreducible (no roots in $\mathbb{Z}$). $\Delta = -4(1)^3 - 27(1)^2 = -31$.
$F = \mathbb{Q}(x_1)$ has degree 3, but the splitting field $L$ has degree 6 ($S_3$), so $F$ is not Galois.

(2) $Gal(L/\mathbb{Q}) = S_3$. $K = \mathbb{Q}(\sqrt{-31})$ is the fixed field of $A_3$.

(3) \textbf{Pedagogical Remark:} The splitting of $p$ is determined by the Frobenius element $\sigma_p$.
(a) If $f$ is not soluble mod $p$, $\sigma_p$ has no fixed points on roots, so $\sigma_p$ is a 3-cycle.
Irreducible cubic mod $p \implies p O_F$ is prime.
(b) $\sigma_p$ is a 3-cycle $\implies$ order 3, so $p O_L$ splits into $6/3 = 2$ primes.
(c) $\sigma_p \in A_3 = Gal(L/K)$, so $\sigma_p|_K = id$.
Thus $p$ splits in $K = \mathbb{Q}(\sqrt{-31})$.
(d) $p$ splits in $\mathbb{Q}(\sqrt{-31}) \iff (-31/p) = 1$.
Since $(31/p) = (-1/p)(-31/p)$, if $(31/p) = -1$ and $p \equiv 3 \pmod 4$, then $(-31/p)=1$.
Actually, the question asks for the implication $(31/p)=-1 \implies$ solvable.
If $f$ is NOT solvable, then $(-31/p)=1$.
For $p \neq 2, 31$, $x^3+x+1$ has a root iff $\sigma_p$ is 1 or a 2-cycle.
$\sigma_p$ is a 3-cycle $\iff (-31/p) = 1$ (even permutation).
Wait, 3-cycles are even. So $(-31/p)=1$ means $\sigma_p \in \{1, (123), (132)\}$.
If $(-31/p)=-1$, then $\sigma_p$ is odd, i.e., a 2-cycle.
A 2-cycle has one fixed point, so $f$ has exactly one root mod $p$.
Thus $(31/p)=-1$ (with some sign adjustments) implies solubility.
\end{solution}

\begin{exercise}
$p$-adic fixed point $\phi$.
\end{exercise}

\begin{solution}
(1) $\phi(x) = x^p + p \sum a_n x^n$.
For $x \equiv y \pmod p$, $x^p \equiv y^p \pmod{p^2}$.
$|\phi(x)-\phi(y)| \approx |x^p-y^p + p \sum \dots| \leq p^{-1} |x-y|$.
(2) Banach Fixed Point Theorem on the compact set $R = a + p\mathbb{Z}_p$.
(3) $F_1$ is constant on disks, $\phi$ maps disk to itself and shrinks it. $\phi^* f = f$ for constant.
For $f(\epsilon_R)=0$, locally constant means $f$ is 0 on a small disk around $\epsilon_R$.
Since $\phi^n$ eventually lands in this disk, $\phi^{*n} f = 0$.
(4) $\phi^* - a$ is invertible on $F_1$ (scalar $1-a$) and on $F_0$ (nilpotent part).
\end{solution}

\begin{exercise}
UFD checks.
\end{exercise}

\begin{solution}
(1) $\mathbb{Z}[\sqrt{6}]$: $h=1$, UFD.
(2) $\mathbb{Z}[\frac{1+\sqrt{-11}}{2}]$: $h=1$, UFD.
(3) $\mathbb{C}[x, y]/(x^2+y^2-1) \cong \mathbb{C}[u, u^{-1}]$, UFD.
(4) $\mathbb{C}[x, y]/(x^3+y^3-1)$: Elliptic curve minus 3 points. Not a UFD (class group of torus is non-trivial).
\end{solution}

\subsection*{Team}
\begin{exercise}
$SL_2(\mathbb{R})$ on $\mathbb{H}$.
\end{exercise}

\begin{solution}
(a) Transitive: $i \to z = x+iy$ via $\begin{pmatrix} \sqrt{y} & x/\sqrt{y} \\ 0 & 1/\sqrt{y} \end{pmatrix}$.
(b) $gi=i \implies ai+b = -c+di \implies a=d, b=-c$. $a^2+c^2=1$. $SO(2)$.
(c) Orbits of $\Gamma(2)$ on $\mathbb{P}^1(\mathbb{Q})$ are the cusps of $X(2)$.
There are 3 orbits: $0, 1, \infty$.
Representative: $0/1, 1/1, 1/0$.
\end{solution}

\begin{exercise}
Product of cosines.
\end{exercise}

\begin{solution}
(a) $C = 1/8$.
(b) $\prod_{k=1}^{(p-1)/2} 2 \cos(k\pi/p) = (2/p)$? No, the formula is:
Let $x^p+1 = (x+1) \prod (x-\zeta^k)$.
Evaluated at $x=1$, $2 = 2 \prod (1-e^{i\dots}) \dots$
Actually $\prod_{k=1}^{(p-1)/2} \cos(n\pi/p) = 2^{-(p-1)/2} (-1)^{(p^2-1)/8}$ is not right.
For $p=7$, $(2/7)=1$, product $1/8$. Correct.
The value is $\frac{1}{2^{(p-1)/2}} \cdot \text{sgn}$, where $\text{sgn}$ is related to $p \pmod 8$.
\end{solution}

\begin{exercise}
$A^{db} = \begin{pmatrix} 0 & I \\ A & 0 \end{pmatrix}$.
\end{exercise}

\begin{solution}
(a) $(A^{db})^2 = A \oplus A$.
$P_{A^{db}}(t) = P_A(t^2)$. $m_{A^{db}}(t) = m_A(t^2)$.
(b) $A \sim B$ iff $A^{db} \sim B^{db}$.
$A^{db} \sim B^{db} \implies A \oplus A \sim B \oplus B$.
In RCF, the invariant factors of $A \oplus A$ are the union of invariant factors of $A$ and $A$.
This uniquely determines the invariant factors of $A$.
So $A \sim B$.
\end{solution}

\begin{exercise}
Groups of order 8.
\end{exercise}

\begin{solution}
Abelian: $C_8, C_4 \times C_2, C_2^3$.
Non-abelian: $D_8$ (symmetries of square), $Q_8$ (quaternions).
\end{solution}

\begin{exercise}
Invariants of $V \otimes V$.
\end{exercise}

\begin{solution}
$V \otimes V \cong \text{End}(V)$. Action is conjugation.
Fixed points are endomorphisms commuting with $O(V)$.
By Schur's Lemma, only scalar multiples of $I$.
Thus $\dim = 1$.
The isomorphism is $v \otimes w \mapsto \phi$ where $\phi(u) = (w, u)v$.
\end{solution}

\begin{exercise}
Norm form $x^3+cy^3+c^2z^3-3cxyz$.
\end{exercise}

\begin{solution}
(a) $(x+\theta y + \theta^2 z)(x+\omega\theta y + \dots)$.
(b) If $c=\theta^3$, $f(x,y,z) = (x+\theta y+\theta^2 z) Q(x,y,z)$.
$f=1 \implies x+\theta y+\theta^2 z = \pm 1$.
This plane has finite intersection with the ellipse $Q=1$? No, $Q=1$ is a cylinder.
Wait, $x^2+\theta^2 y^2 + \dots$ is not definite if $\theta < 0$.
But for $\theta \in \mathbb{Z}$, it defines a family of points.
Actually, if $c$ is a cube, the field is $\mathbb{Q}$, so only finitely many solutions (bounded by $Q=1$ which is an ellipse in the plane).
(c) If $c$ not a cube, it's the norm form of $\mathbb{Q}(\sqrt[3]{c})$.
The unit group has rank 1 ($r_1=1, r_2=1$).
Infinite solutions by Dirichlet's Unit Theorem.
\end{solution}
